\documentclass[12pt]{article}
\usepackage{amsmath, amssymb, geometry}
\geometry{a4paper, margin=1in}
\setcounter{MaxMatrixCols}{100}

\begin{document}

\title{Phase Space Debug Report: \texttt{m036}}
\author{Automated Pipeline}
\date{\today}
\maketitle

\section{Gluing Equations (SnapPy)}
Manifold: \texttt{m036},\ $n = 4$ tetrahedra,\ $r = 1$ cusp.
Columns ordered as $Z_1, Z_1', Z_1'',\, Z_2, Z_2', Z_2'',\, Z_3, Z_3', Z_3'',\, Z_4, Z_4', Z_4''$.

\textbf{Edge equations ($C$ matrix, shape $4 \times 12$):}
\[ C = \begin{bmatrix}
2 & 0 & 1 & 1 & 2 & 0 & 1 & 1 & 0 & 1 & 1 & 0 \\
0 & 1 & 0 & 0 & 0 & 1 & 1 & 0 & 0 & 1 & 0 & 0 \\
0 & 1 & 1 & 0 & 0 & 0 & 0 & 1 & 1 & 0 & 0 & 1 \\
0 & 0 & 0 & 1 & 0 & 1 & 0 & 0 & 1 & 0 & 1 & 1
\end{bmatrix} \]
Note: $n - r = 3$ edge equations are independent; the fourth is their sum
(all four are shown as SnaPy returns them).

\textbf{Cusp equations (row order: $\mu_0,\, \lambda_0$):}
\[ C_{\mathrm{cusp}} = \begin{bmatrix}
 0 &  0 &  1 &  0 &  2 &  0 &  0 &  0 & -1 &  0 &  0 &  0 \\
-2 &  0 &  0 &  0 &  2 &  0 &  0 &  0 &  0 &  0 &  0 &  0
\end{bmatrix} \]

\section{Geometric Basis}
The pipeline finds $n - r = 3$ independent internal edges.
These are classified into \textbf{2 hard} edges (carrying fugacity variables $\eta_1, \eta_2$)
and \textbf{1 easy} edge (whose charge is summed freely, with no fugacity variable).

\textbf{Easy edge criterion:} a basis edge whose 3n-vector has exactly one non-zero
entry per tetrahedron, lying in the $Z'$ slot.  Such edges contribute only
$I_\Delta(0,e)$ factors and can be freely summed.

\begin{itemize}
    \item \textbf{Easy Edge $Q_3$:}\ Vector $= [0,1,0,\; 0,0,1,\; 1,0,0,\; 1,0,0]$
    \item \textbf{Hard Edge $Q_1$:}\ Vector $= [2,0,1,\; 1,2,0,\; 1,1,0,\; 1,1,0]$
    \item \textbf{Hard Edge $Q_2$:}\ Vector $= [0,0,0,\; 1,0,1,\; 0,0,1,\; 0,1,1]$
\end{itemize}

Since there are 2 hard edges, the refined index carries \textbf{two fugacity variables}
$\eta_1, \eta_2$.

\section{Neumann-Zagier Matrix}
Row structure of $g_{NZ} \in \mathrm{Sp}(8,\mathbb{Q})$
($2n \times 2n$ with $n=4$, columns in block order $Z_1, Z_2, Z_3, Z_4,\, Z_1'', Z_2'', Z_3'', Z_4''$):
\[
\text{rows } 0: \mu_0 \quad
\text{rows } 1,2: Q_1,\, Q_2\ (\text{hard}) \quad
\text{row } 3: Q_3\ (\text{easy}) \quad
\text{row } 4: \tfrac{\lambda_0}{2} \quad
\text{rows } 5,6,7: \Gamma_1,\, \Gamma_2,\, \Gamma_3
\]

\textbf{$g_{NZ}$ Matrix:}
\[ g_{NZ} = \begin{bmatrix}
 0 & -2 &  0 &  0 &  1 & -2 & -1 &  0 \\
 2 & -1 &  0 &  0 &  1 & -2 & -1 & -1 \\
 0 &  1 &  0 & -1 &  0 &  1 &  1 &  0 \\
-1 &  0 &  1 &  1 & -1 &  1 &  0 &  0 \\
-1 & -1 &  0 &  0 &  0 & -1 &  0 &  0 \\
 0 &  0 &  0 &  1 &  0 &  0 &  0 &  0 \\
 0 & -2 & -2 &  1 &  1 & -3 & -1 &  0 \\
 0 & -2 &  0 &  2 &  0 & -2 & -1 &  0
\end{bmatrix} \]

\textbf{Affine Shift $\nu$:}
\[
\nu_X = \begin{bmatrix} 2 \\ 2 \\ -1 \\ -1 \end{bmatrix}, \qquad
\nu_P = \begin{bmatrix} 1 \\ 0 \\ 0 \\ 0 \end{bmatrix}
\]
The rows of $\nu_X$ correspond to position rows (rows $0$--$3$: $\mu_0, Q_1, Q_2, Q_3$)
and $\nu_P$ to momentum rows (rows $4$--$7$: $\lambda_0/2, \Gamma_1, \Gamma_2, \Gamma_3$).

\section{Inversion and Local Charges}
\textbf{Inverse Matrix $g_{NZ}^{-1}$ (all-integer, verified $g_{NZ} g_{NZ}^{-1} = I$):}
\[ g_{NZ}^{-1} = \begin{bmatrix}
 0 &  0 &  1 &  0 & -1 & -1 &  0 &  1 \\
-1 &  0 & -3 & -2 &  2 &  2 & -1 & -1 \\
 0 &  0 & -1 & -1 &  1 &  1 & -1 &  0 \\
 0 &  0 &  0 &  0 &  0 &  1 &  0 &  0 \\
 1 &  0 &  0 &  0 &  0 &  2 &  0 & -1 \\
 1 &  0 &  2 &  2 & -2 & -1 &  1 &  0 \\
 0 &  0 &  2 &  0 &  0 &  0 &  0 &  1 \\
 0 & -1 & -1 & -2 &  0 &  0 & -1 &  1
\end{bmatrix} \]

The master charge vector is
\[
\kappa = \bigl[m,\; 0,\; 0,\; 0;\; e,\; e_{h_1},\; e_{h_2},\; e_{\mathrm{easy}}\bigr]^T,
\]
where $m$ is the cusp meridian charge, the three zeros are internal position charges
(always fixed to zero in the summation), $e$ is the cusp longitude$/2$ charge,
$e_{h_1}, e_{h_2} \in \tfrac{1}{2}\mathbb{Z}$ are the hard-edge momentum charges
(carrying fugacities $\eta_1^{2e_{h_1}}, \eta_2^{2e_{h_2}}$),
and $e_{\mathrm{easy}} \in \tfrac{1}{2}\mathbb{Z}$ is the easy-edge momentum charge
(summed freely with no fugacity).

The local tetrahedron charges $\Delta = g_{NZ}^{-1}\,\kappa$:
\begin{align*}
\Delta_1 &: \quad
  \tilde{m}_1 = -e - e_{h_1} + e_{\mathrm{easy}},
  \qquad
  \tilde{e}_1 = m + 2e_{h_1} - e_{\mathrm{easy}} \\[4pt]
\Delta_2 &: \quad
  \tilde{m}_2 = -m + 2e + 2e_{h_1} - e_{h_2} - e_{\mathrm{easy}},
  \qquad
  \tilde{e}_2 = m - 2e - e_{h_1} + e_{h_2} \\[4pt]
\Delta_3 &: \quad
  \tilde{m}_3 = e + e_{h_1} - e_{h_2},
  \qquad
  \tilde{e}_3 = e_{\mathrm{easy}} \\[4pt]
\Delta_4 &: \quad
  \tilde{m}_4 = e_{h_1},
  \qquad
  \tilde{e}_4 = -e_{h_2} + e_{\mathrm{easy}}
\end{align*}

The topological framing phase power
$\phi = \vec{m}_{\mathrm{full}} \cdot \nu_P - \vec{e}_{\mathrm{full}} \cdot \nu_X$,
where $\vec{m}_{\mathrm{full}} = (m, 0, 0, 0)$ and
$\vec{e}_{\mathrm{full}} = (e,\, e_{h_1},\, e_{h_2},\, e_{\mathrm{easy}})$.
Substituting $\nu_P = (1,0,0,0)$ and $\nu_X = (2,2,-1,-1)$:
\[
\phi
= m \cdot 1 - \bigl(e \cdot 2 + e_{h_1} \cdot 2 + e_{h_2} \cdot (-1)
  + e_{\mathrm{easy}} \cdot (-1)\bigr)
= m - 2e - 2e_{h_1} + e_{h_2} + e_{\mathrm{easy}}.
\]

\section{Final Summation Formula}
The refined index is:
\[
I_{m036}(m, e)
= \sum_{\substack{e_{h_1},\, e_{h_2},\, e_{\mathrm{easy}} \\ \in\, \frac{1}{2}\mathbb{Z}}}
   \eta_1^{2e_{h_1}}\,\eta_2^{2e_{h_2}}
   \cdot (-q^{1/2})^{\phi}
   \cdot \prod_{i=1}^{4} I_{\Delta}(\tilde{m}_i,\, \tilde{e}_i),
\]
where $\phi = m - 2e - 2e_{h_1} + e_{h_2} + e_{\mathrm{easy}}$
and the $\tilde{m}_i, \tilde{e}_i$ are given above.

\textbf{Evaluated output at $(m, e) = (0, 0)$ (truncated at $q^4$):}
\begin{align*}
I &= 1 \\
&\quad + q\!\left(-\eta_2^{-1} - 3 - \eta_1\eta_2\right) \\
&\quad + q^2\!\left(\eta_1^{-1}\eta_2^{-2} + \eta_1^{-1}\eta_2^{-1}
                   - \eta_2^{-1} - 1 + \eta_2 - \eta_1\eta_2 + \eta_1\eta_2^{2}\right) \\
&\quad + q^3\!\left(\eta_1^{-1}\eta_2^{-3} + 3\eta_1^{-1}\eta_2^{-2} + \eta_2^{-2}
                   + 3\eta_2^{-1} + 10 + \eta_1
                   + 3\eta_1\eta_2 + 3\eta_1\eta_2^{2}
                   + \eta_1^{2}\eta_2^{2} + \eta_1^{2}\eta_2^{3}\right) \\
&\quad + q^4\!\left(\eta_1^{-1}\eta_2^{-3} + \eta_1^{-1}\eta_2^{-2}
                   - 4\eta_1^{-1}\eta_2^{-1} + \eta_1^{-1}
                   + \eta_2^{-2} + 10\eta_2^{-1} + 14 - 4\eta_2
                   + 3\eta_1 + 10\eta_1\eta_2 + \eta_1\eta_2^{2}
                   + \eta_1^{2}\eta_2^{2} + \eta_1^{2}\eta_2^{3}\right)
\end{align*}

\noindent\textbf{Specialisation to 3D index} (setting $\eta_1 = \eta_2 = 1$,
projecting over all fugacity sectors):
\[
I_{3D}(q) = 1 - 5q + q^2 + 27q^3 + 34q^4 + \cdots
\]
(This matches the direct 3D-index computation via \texttt{compute\_index\_3d\_python}.)

\bigskip
\noindent\textbf{Raw key convention} (as stored in \texttt{RefinedIndexResult}):
a monomial $q^k\,\eta_1^a\,\eta_2^b$ is stored under key $(2k,\; 2a,\; 2b)$,
i.e.\ all exponents doubled to allow half-integer $q$-powers.

\bigskip
\noindent\textbf{Integrality / half-integrality of summation variables.}
The charges $e_{h_1}, e_{h_2}, e_{\mathrm{easy}}$ are summed over $\tfrac{1}{2}\mathbb{Z}$.
For a given $(m, e)$, only those values that make every local charge
$\tilde{m}_i, \tilde{e}_i$ simultaneously a half-integer (relative to the
SnaPy orientation) contribute.  The pipeline enforces this by requiring
\[
g_{NZ}^{-1}\,\kappa \in \left(\tfrac{1}{2}\mathbb{Z}\right)^{2n}.
\]

\end{document}
